\phantomsection\label{fig:vol01:warring-qin:chu-eastward-court}
\begin{center}
\begin{tikzpicture}[x=.9cm,y=.9cm,font=\small]
  \fill[paper] (0,0) rectangle (10.8,5.7);
  \draw[blue!55,line width=1pt] (.45,4.65) .. controls (2.6,4.15) and (4.4,4.75) .. (6.35,4.15);
  \node[text=blue!70!black,above] at (3.55,4.43) {汉水方向};
  \draw[blue!50,line width=.9pt] (1.1,1.0) .. controls (3.35,1.95) and (5.45,1.15) .. (9.9,.7);
  \node[text=blue!70!black,below] at (6.55,1.18) {淮水方向};
  \draw[gray!55,line width=1pt] (1.0,5.15) .. controls (1.85,3.85) and (1.4,2.0) .. (2.35,.5);
  \node[rotate=76,text=gray!70] at (1.62,3.5) {西部山地};

  \fill[purple!17] (1.55,1.65) .. controls (2.7,1.45) and (4.4,1.7) .. (5.05,2.75)
    .. controls (4.6,3.75) and (3.2,4.05) .. (1.85,3.45) -- cycle;
  \node[align=center] at (3.18,3.0) {楚\\旧都腹地};
  \fill[dynasty] (3.0,2.6) circle (.14);
  \node[above,align=center] at (3.0,2.67) {郢\\前278年失守};

  \fill[timeline] (6.0,2.85) circle (.14);
  \node[above,align=center] at (6.0,2.92) {陈\\王廷东迁};
  \fill[timeline!23] (7.25,1.75) ellipse (.95 and .65);
  \fill[timeline] (7.25,1.75) circle (.14);
  \node[right,align=left] at (7.4,1.75) {寿春\\后期中心};

  \draw[-{Stealth[length=2mm]},ultra thick,color=dynasty] (1.38,3.75) -- (2.78,2.72);
  \node[align=center,text=dynasty] at (1.55,2.65) {秦军沿\\上游东进};
  \draw[-{Stealth[length=2mm]},ultra thick,dashed,color=timeline] (3.2,2.55) .. controls (4.45,3.75) .. (5.82,2.98);
  \draw[-{Stealth[length=2mm]},thick,dashed,color=timeline] (6.15,2.7) -- (7.13,1.88);
  \node[align=center,text=timeline] at (4.68,3.55) {王廷转移方向\\路线未详};
\end{tikzpicture}

\smallskip
\small\textit{图\quad 楚都东迁与南方纵深示意：据《史记·白起王翦列传》《楚世家》与《资治通鉴》提示前278年郢失守后楚廷向陈、寿春方向重组的空间压力。虚线只表示政治中心转移的约略方向，并非人口迁徙规模或确切行程；疆界、河道和道路均为约略示意。}
\end{center}
